%% 由 build/emit_tex.py 生成（生成物，勿手改）；定制请放 user_style.tex / user_meta.yaml
\chapter{概率统计基础}\label{ux6982ux7387ux7edfux8ba1ux57faux7840}

\begin{quote}
一句话：概率给''不确定性''建模，统计从数据里''猜''出规律。本附录覆盖描述统计、常用分布、推断（检验/置信区间）、回归与时间序列。
\end{quote}

\section{C.1
直觉与描述统计}\label{c.1-ux76f4ux89c9ux4e0eux63cfux8ff0ux7edfux8ba1}

\begin{itemize}
\tightlist
\item
  \textbf{随机变量}：取值带概率的变量。期望 \texttt{E{[}X{]}} =
  中心位置；方差 \texttt{Var(X)\ =\ E{[}(X-E{[}X{]})²{]}} = 分散程度。
\item
  \textbf{样本统计量}：从数据算出来的''估计''------均值、中位数、标准差、分位数、相关系数。
\item
  \textbf{为什么先看分布}：均值会被极端值拉偏；分布形状（偏斜/厚尾）决定该用哪种方法。
\end{itemize}

常用量速查：

{\def\LTcaptype{none} % do not increment counter
\begin{longtable}[]{@{}lll@{}}
\toprule\noalign{}
量 & 定义 & 说明 \\
\midrule\noalign{}
\endhead
\bottomrule\noalign{}
\endlastfoot
均值 \texttt{μ\^{}\ =\ (1/n)Σx\_i} & 集中趋势 & 受异常值影响 \\
中位数 & 排序中间值 & 稳健 \\
方差/标准差 & \texttt{s²\ =\ Σ(x\_i\ -\ x̄)²/(n-1)} & 离散程度（样本用
n-1） \\
协方差/相关 & \texttt{Cov(X,Y)}、\texttt{ρ\ =\ Cov/(σx\ σy)} &
线性相关方向/强度 \\
分位数 & 排序后百分比位置 & Q1/Q3/IQR（箱线图） \\
\end{longtable}
}

\section{C.2
常用分布速查}\label{c.2-ux5e38ux7528ux5206ux5e03ux901fux67e5}

{\def\LTcaptype{none} % do not increment counter
\begin{longtable}[]{@{}lll@{}}
\toprule\noalign{}
分布 & 记号 & 用在哪 \\
\midrule\noalign{}
\endhead
\bottomrule\noalign{}
\endlastfoot
正态 & \texttt{N(μ,\ σ²)} & 近似很多现象；中心极限定理的主角 \\
二项 & \texttt{B(n,\ p)} & 成功/失败次数 \\
泊松 & \texttt{Poisson(λ)} & 一段时间内事件数 \\
均匀 & \texttt{U(a,\ b)} & 随机数/先验 \\
t / F / χ² & \texttt{t(n)}, \texttt{F}, \texttt{χ²} &
小样本推断/方差分析 \\
\end{longtable}
}

\begin{quote}
\textbf{中心极限定理}：大量独立同分布随机变量的均值，近似正态分布------这是很多统计方法（包括回归的置信区间）的底气。
\end{quote}

\section{C.3
统计推断：检验与置信区间}\label{c.3-ux7edfux8ba1ux63a8ux65adux68c0ux9a8cux4e0eux7f6eux4fe1ux533aux95f4}

\begin{itemize}
\tightlist
\item
  \textbf{参数估计}：点估计（均值）→
  区间估计（置信区间：\texttt{x̄\ ±\ t\_\{α/2\}·s/√n}）。
\item
  \textbf{假设检验}：先立原假设 \texttt{H0}（如''均值=0''），看数据在
  \texttt{H0} 下多''极端''（p 值）；\texttt{p\ \textless{}\ 0.05} → 拒绝
  \texttt{H0}。
\item
  \textbf{两类错误}：第一类（误拒，α）、第二类（漏拒，β）；功效 = 1-β。
\item
  \textbf{常见检验}：t
  检验（均值）、卡方（拟合/列联）、F/ANOVA（多组均值）、KS（分布）。
\item
  \textbf{重要提醒}：p 值小 ≠ 实际重要；报告效应量 + 置信区间。
\end{itemize}

\section{C.4
回归与方差分析（建模视角）}\label{c.4-ux56deux5f52ux4e0eux65b9ux5deeux5206ux6790ux5efaux6a21ux89c6ux89d2}

\begin{itemize}
\tightlist
\item
  \textbf{线性回归}：\texttt{y\ =\ β0\ +\ β1x1\ +\ …\ +\ ε}；OLS
  最小化残差平方和（附录 A 最小二乘）。
\item
  \textbf{假设}：线性、独立、同方差、正态（对推断）；违背时用稳健标准误/变换。
\item
  \textbf{\texttt{R²}}：解释变量能解释的变异比例（\texttt{R²\ =\ 1\ -\ SSE/SST}），过大要提防过拟合。
\item
  \textbf{ANOVA}：检验多个组的均值是否相等（\texttt{F}
  检验，组间方差/组内方差）。
\item
  \textbf{逻辑回归}：把线性预测过 sigmoid
  映射成概率：\texttt{P(y=1)\ =\ σ(β\^{}T\ x)}；损失用交叉熵（与附录 E
  衔接）。
\end{itemize}

\section{C.5
时间序列基础}\label{c.5-ux65f6ux95f4ux5e8fux5217ux57faux7840}

{\def\LTcaptype{none} % do not increment counter
\begin{longtable}[]{@{}lll@{}}
\toprule\noalign{}
概念 & 直觉 & 常用工具 \\
\midrule\noalign{}
\endhead
\bottomrule\noalign{}
\endlastfoot
平稳性 & 统计特性不随时间漂移 & 单位根检验（ADF） \\
自相关/ACF & 与过去的相关 & 确定滞后阶数 \\
差分 & \texttt{y\_t\ -\ y\_\{t-1\}} 消除趋势 & 转平稳 \\
移动平均/平滑 & 去噪/看趋势 & \texttt{rolling} \\
分解 & 趋势 + 季节 + 残差 & \texttt{seasonal\_decompose} \\
ARIMA & AR（自回归）+ I（差分）+ MA（移动平均误差） & 预测/理解 \\
\end{longtable}
}

\section{C.6 Python 对应}\label{c.6-python-ux5bf9ux5e94}

\begin{Shaded}
\begin{Highlighting}[]
\ImportTok{import}\NormalTok{ numpy }\ImportTok{as}\NormalTok{ np, pandas }\ImportTok{as}\NormalTok{ pd, statsmodels.api }\ImportTok{as}\NormalTok{ sm}
\ImportTok{from}\NormalTok{ scipy }\ImportTok{import}\NormalTok{ stats}

\NormalTok{x }\OperatorTok{=}\NormalTok{ np.array([}\DecValTok{68}\NormalTok{, }\DecValTok{91}\NormalTok{, }\DecValTok{55}\NormalTok{, }\DecValTok{76}\NormalTok{, }\DecValTok{88}\NormalTok{])}
\BuiltInTok{print}\NormalTok{(np.mean(x), np.std(x, ddof}\OperatorTok{=}\DecValTok{1}\NormalTok{), np.percentile(x, }\DecValTok{75}\NormalTok{))}

\CommentTok{\# 单样本 t 检验}
\NormalTok{stats.ttest\_1samp(x, }\DecValTok{60}\NormalTok{)}

\CommentTok{\# 线性回归}
\NormalTok{X }\OperatorTok{=}\NormalTok{ sm.add\_constant(np.array([}\DecValTok{1}\NormalTok{, }\DecValTok{2}\NormalTok{, }\DecValTok{3}\NormalTok{, }\DecValTok{4}\NormalTok{, }\FloatTok{5.}\NormalTok{]))}
\NormalTok{model }\OperatorTok{=}\NormalTok{ sm.OLS(np.array([}\DecValTok{2}\NormalTok{, }\DecValTok{4}\NormalTok{, }\DecValTok{5}\NormalTok{, }\DecValTok{8}\NormalTok{, }\DecValTok{9}\NormalTok{]), X).fit()}
\BuiltInTok{print}\NormalTok{(model.summary())   }\CommentTok{\# 看 coef / p / R²}

\CommentTok{\# 时间序列分解}
\ImportTok{from}\NormalTok{ statsmodels.tsa.seasonal }\ImportTok{import}\NormalTok{ seasonal\_decompose}
\NormalTok{df }\OperatorTok{=}\NormalTok{ pd.read\_csv(}\StringTok{"data.csv"}\NormalTok{, index\_col}\OperatorTok{=}\DecValTok{0}\NormalTok{, parse\_dates}\OperatorTok{=}\VariableTok{True}\NormalTok{)}
\NormalTok{seasonal\_decompose(df[}\StringTok{"sales"}\NormalTok{], model}\OperatorTok{=}\StringTok{"additive"}\NormalTok{).plot()}
\end{Highlighting}
\end{Shaded}

\section{C.7 常见误区}\label{c.7-ux5e38ux89c1ux8befux533a}

{\def\LTcaptype{none} % do not increment counter
\begin{longtable}[]{@{}
  >{\raggedright\arraybackslash}p{(\linewidth - 2\tabcolsep) * \real{0.5000}}
  >{\raggedright\arraybackslash}p{(\linewidth - 2\tabcolsep) * \real{0.5000}}@{}}
\toprule\noalign{}
\begin{minipage}[b]{\linewidth}\raggedright
误区
\end{minipage} & \begin{minipage}[b]{\linewidth}\raggedright
正确
\end{minipage} \\
\midrule\noalign{}
\endhead
\bottomrule\noalign{}
\endlastfoot
p 值小 = 影响大 & p 值是''在 H0
下看到这个数据的概率''，与效应量是两回事 \\
忽视样本量 & 样本大时微小差异也显著；看置信区间宽度 \\
用均值代替看分布 & 先画直方图/箱线图 \\
非平稳直接回归 & 时序先检验平稳性/差分 \\
把相关系数当因果 & 相关 ≠ 因果 \\
\end{longtable}
}

\section{C.8
使用章节与下游衔接}\label{c.8-ux4f7fux7528ux7ae0ux8282ux4e0eux4e0bux6e38ux8854ux63a5}

\begin{itemize}
\tightlist
\item
  章节：0 前置基础（成绩统计）、4 Pandas（分组聚合/描述统计）、7
  Statsmodels（回归/ANOVA/时序）、3 SciPy（假设检验）；
\item
  下游：intro-mathmodel 第 7 章（权重生成与评价模型）、第 8
  章（时间序列与投资模型）、第 9 章（机器学习与统计模型）；
\item
  参考：statsmodels 官方文档、ThinkStats2、中文《概率论与数理统计》（见
  references.md）。
\end{itemize}
